\subsection{Gluing, resolution, and the information in groupoids}

\begin{saresult}{A6.001}{Proved}{The six functors in the gluing category}
Let $\mathcal A,\mathcal B$ be abelian categories and let
$T:\mathcal B\longrightarrow\mathcal A$ be an additive left exact functor.
Write $\mathcal C_T$ for the category of triples
$(F,G,u)$ with $F\in\mathcal A$, $G\in\mathcal B$, and $u:F\to TG$.
A morphism $(a,b)$ satisfies $T(b)u=u'a$.
Then $\mathcal C_T$ is abelian and has the recollement functors
\[
\begin{array}{lll}
i^*(F,G,u)=F,&i_*F=(F,0,0),&i^!(F,G,u)=\ker u,\\
j_!G=(0,G,0),&j^*(F,G,u)=G,&j_*G=(TG,G,1).
\end{array}
\]
Their adjunctions are $i^*\dashv i_*\dashv i^!$ and
$j_!\dashv j^*\dashv j_*$. Both $i_*$ and $j_!,j_*$ are fully faithful,
and $\ker j^*=\operatorname{im}i_*$. The functors $i^*,i_*,j_!,j^*$
are exact; $i^!,j_*$ are left exact. There are natural exact sequences
\[
0\to j_!j^*C\to C\to i_*i^*C\to0,
\qquad
0\to i_*i^!C\to C\to j_*j^*C
 \to i_*\operatorname{coker}u\to0.
\]
For the displayed triples on 121/039, take
$\mathcal A=\operatorname{Sh}_\Lambda(Y)$,
$\mathcal B=\operatorname{Sh}_\Lambda(X)$, and $T=p_*i^{-1}$.
The adjunction $p^{-1}\dashv p_*$ identifies $u$ with the manuscript's
$\alpha:p^{-1}F\to i^{-1}G$. Here $i^{-1}$ means ordinary sheaf
restriction; it is not the extraordinary restriction $i^!$.
\end{saresult}
\begin{proof}
The kernel of $(a,b)$ has components $\ker a,\ker b$ because
$T\ker b=\ker Tb$. Its structure arrow is induced by $u$.
The cokernel has components $\operatorname{coker}a,
\operatorname{coker}b$: the composite
$F'\xrightarrow{u'}TG'\to T\operatorname{coker}b$ kills $a(F)$
and therefore factors through $\operatorname{coker}a$.
Consequently the canonical coimage--image arrow is an isomorphism
on both components, and hence is an isomorphism in $\mathcal C_T$.
This proves abelianness and the four asserted exactness statements.

A morphism $C\to i_*A$ is an arbitrary arrow $F\to A$.
A morphism $i_*A\to C$ is an arrow $A\to F$ killed by $u$, and
therefore an arrow $A\to\ker u$. A morphism $j_!B\to C$ is an
arbitrary arrow $B\to G$. Finally, a morphism $C\to j_*B$ is
uniquely specified by $b:G\to B$, since its first component must
be $T(b)u$. These four calculations prove all the adjunctions.
The units or counits on the outer categories are identities, which
proves full faithfulness. An object with $G=0$ is precisely $i_*F$.
The first exact sequence has component maps $(0,1)$ and $(1,0)$.
The middle arrow in the second is $(u,1)$; its kernel and cokernel
are exactly the displayed ones. Left exactness of the two right
adjoints also follows directly from these formulas.

The functor $i^{-1}$ for sheaves of modules over a fixed ring is exact,
and $p_*$ is left exact. Thus this application meets every hypothesis.
It recovers the manuscript's $i'_*,i'^*,p'_*,p'^*$ as
$i_*,i^*,j_*,j^*$ in the abstract diagram, and supplies its missing
right adjoint and extension by zero. These names in the abstract
diagram do not by themselves identify $p'$ with an open immersion
of the topological pushout; that separate question is A6.002.
\end{proof}

\begin{saresult}{A6.002}{Proved under stated hypotheses}{Actual sheaf gluing and the necessary correction on page 039}
Let $i:Y\hookrightarrow Z$ be a closed embedding and
$j:U=Z\setminus Y\hookrightarrow Z$ its open complement. For sheaves
of modules over any fixed ring $\Lambda$, restriction induces an
equivalence
\[
\operatorname{Sh}_\Lambda(Z)\simeq
\mathcal C_{i^{-1}j_*}.
\]
It identifies all six functors with A6.001. In particular the right
adjoint to restriction to the open part is $j_*$, and the right
adjoint to extension from the closed part is the kernel of the
specialization arrow.

For a boundary pushout $Z=Y\cup_pX$ with $p:\partial X\to Y$,
the unrestricted triples of page 039, with $G$ an arbitrary sheaf on
the whole resolved $X$, are in general a larger category than
$\operatorname{Sh}_\Lambda(Z)$. They become the usual gluing formula
when $G$ denotes a sheaf on $U=X\setminus\partial X$ and the gluing
functor is $i^{-1}j_*$. If a compatible mapping-cylinder collar is
chosen and $G$ is pulled back from a sheaf $H$ on the link in that
collar, this functor is $p_*H$ when the radial interval is the ordinary
real interval. This last radial simplification is not asserted for
all ordinary topological sheaves on a non-Archimedean order topology.
The open--closed formula and the six abstract adjoints above remain
valid for ordinary sheaves on every topological space.
\end{saresult}
\begin{proof}
For a triple $(F,G,u:F\to i^{-1}j_*G)$ form the fiber product of
sheaves
\[
 E=i_*F\mathbin{\times}_{i_*i^{-1}j_*G}j_*G.
\]
The two arrows to the base are $i_*u$ and the restriction unit.
The exact restriction $i^{-1}$ preserves this finite limit and
identifies its restriction with $F\times_{i^{-1}j_*G}i^{-1}j_*G=F$.
Restriction to $U$ identifies it with $G$. Its specialization arrow
is $u$. Conversely, the natural map from any sheaf $E$ to this fiber
product for $(i^{-1}E,j^{-1}E,u_E)$ is an isomorphism: on a stalk in
$Y$ it is the identity of $E_y$, and on a stalk in $U$ it is the
identity of $E_x$. A sheaf morphism inducing an isomorphism at every
stalk is an isomorphism. This proves both inverse constructions,
including their action on morphisms. A6.001 then computes the
adjoints. The elementary properties of restriction and closed
extension used here are also stated in the opened Stacks Project,
Sections 6.31--6.32, tags 009Z and 00AD~\cite{StacksSheavesOpenClosed}.

For the correction, take $X=[0,\infty)$, $Y=\partial X=\{0\}$,
and $p=1$. The pushout is $X$ again. Choose a nonzero $\Lambda$-module
$K$ and the triple $(0,i_*K,0)$ in the page-039 category.
It is nonzero, but its proposed restriction to the closed stratum
and its restriction to the open stratum are both zero. No nonzero
sheaf on $X$ has both restrictions zero, by the stalk criterion.
Thus no equivalence with the asserted geometric restrictions can
hold for these unrestricted triples. This counterexample uses a
genuine collared boundary and the identity boundary fibration.

Finally, in a locally trivial mapping-cylinder chart above an open
set $V\subset Y$, punctured collar neighborhoods are
$p^{-1}V\times(0,\epsilon)$ with the endpoints at radius zero
identified by $p$. A radially pulled-back sheaf has sections there
equal to $H(p^{-1}V)$, compatibly with restriction in $V$ and in
$\epsilon$: a locally constant section along a connected real
interval is constant. Taking the neighborhood germs along $Y$
therefore gives
$p_*H$. Without radial constancy a single level of a collar need
not compute these germs, which is why that additional hypothesis
was stated explicitly.
\end{proof}

\begin{saresult}{A6.003}{Proved}{Derived gluing with its six functors}
For stable categories $\mathcal A,\mathcal B$ and an exact functor
$T:\mathcal B\to\mathcal A$, let $\mathcal C_T$ be the homotopy
gluing category of arrows $F\to TG$. Its six functors have the same
formulas as A6.001, with $\ker u$ replaced by $\operatorname{fib}u$.
The exact sequences become the distinguished triangles
\[
j_!j^*C\longrightarrow C\longrightarrow i_*i^*C\longrightarrow,
\qquad
i_*i^!C\longrightarrow C\longrightarrow j_*j^*C\longrightarrow.
\]
For derived sheaves in the preceding open--closed situation the
gluing functor is $T=i^{-1}Rj_*$, so the extraordinary restriction
is the homotopy fiber of $F\to i^{-1}Rj_*G$.
\end{saresult}
\begin{proof}
The mapping space between two arrows is the homotopy pullback
of $\operatorname{Map}_{\mathcal A}(F,F')$ and
$\operatorname{Map}_{\mathcal B}(G,G')$ over
$\operatorname{Map}_{\mathcal A}(F,TG')$.
Substituting the three inclusions gives exactly the four mapping
space adjunctions of A6.001. For $i^!$ the resulting mapping space
is the fiber of
$\operatorname{Map}(A,F)\to\operatorname{Map}(A,TG)$, hence is
$\operatorname{Map}(A,\operatorname{fib}u)$.
Fibers and cofibers are computed componentwise because $T$ is exact.
This proves the two triangles. In derived sheaves the inverse
construction is the homotopy fiber product
$i_*F\times^h_{i_*i^{-1}Rj_*G}Rj_*G$.
Restriction to $Y$ and $U$ identifies it with $F$ and $G$ exactly
as before; these restrictions jointly detect equivalences through
cohomology stalks. It follows that the construction and restriction
are inverse equivalences. This argument uses the derived gluing
category; it does not silently identify the derived category of an
arbitrary comma category with a comma category of ordinary complexes.
\end{proof}

\begin{saresult}{A6.004}{Proved under stated hypotheses}{The precise finite resolution datum}
The following is a precise form of the data on 121/041--047.
Fix a finite partially ordered set $I$. For each $i$ give a compact
manifold with embedded corners $M_i$ of dimension $d_i$. For each
$j<i$ give a possibly empty collective boundary hypersurface
$H_{ij}\subset M_i$ and a locally trivial fibration
$p_{ij}:H_{ij}\to M_j$ of manifolds with corners, whenever
$H_{ij}\ne\varnothing$. A collective hypersurface is a disjoint
union of boundary hypersurfaces with a specified common base.
Require the following conditions.
\begin{enumerate}
\item Every boundary hypersurface occurs in exactly one $H_{ij}$.
An intersection of several nonempty $H_{ij}$ can occur only when
their lower labels form a chain. The intersection is a corner of
the expected codimension.
\item For $j<k<i$ one has
\[
H_{ij}\cap H_{ik}=p_{ik}^{-1}(H_{kj}),\qquad
p_{ij}=p_{kj}p_{ik}\quad\hbox{on }H_{ij}\cap H_{ik}.
\]
Restrictions of the fibrations to these corners are locally trivial.
Their entire finite diagram over a sufficiently small open set in
$M_j$ has a simultaneous local trivialization. This last condition
specifies compatibility as a family, and not merely the separate
local triviality of each arrow.
\item Compatible collars are part of the datum. They commute on
corners and preserve all the displayed maps: a collar for a lower
face $H_{ij}$, restricted to $H_{ik}$ with $j<k$, projects by
$p_{ik}$ to the corresponding collar of $H_{kj}$ in $M_k$;
a collar for $H_{ik}$, restricted to $H_{ij}$, is vertical for
$p_{ij}$. Their products give corner charts with independent
nonnegative radius coordinates.
\end{enumerate}
Here and below an empty face contributes no condition. The fiber of
$p_{ij}$ has dimension $d_i-d_j-1$. Its boundary consists exactly of
the corners with labels $j<k<i$; the part of $\partial H_{ij}$ over
$\partial M_j$ consists of the labels $k<j$. Thus
\[
\partial H_{ij}=
\bigcup_{k<j}p_{ij}^{-1}(H_{jk})
\ \cup\ \bigcup_{j<k<i}(H_{ij}\cap H_{ik}).
\]
In particular a fiber has empty boundary when $j$ is an immediate
predecessor of $i$. For a chain $i>j_1>\cdots>j_r$ the total corner
$H_{ij_1}\cap\cdots\cap H_{ij_r}$ has dimension $d_i-r$.
The subscripts in the dimension formula refer to $i$, not to $j_1$.
\end{saresult}
\begin{proof}
In a product chart for a fibration of corner manifolds, the boundary
is the union of the product over the boundary of the base and the
product with the boundary of the fiber. At a point of $H_{ij}$ all
other incident faces have labels comparable to $j$ by condition 1.
If $k<j$, condition 2 identifies the corner with the inverse image
of $H_{jk}$, giving the base-boundary part. If $j<k<i$, the same
condition identifies the corner as a fiber-boundary face over $M_j$.
Condition 1 exhausts the faces, so there is no third contribution.
Subtracting dimensions in a local product chart gives
$(d_i-1)-d_j$ for the fiber. Each independent collar coordinate
set to zero lowers the total dimension by one, proving the final
formula for every chain length.

These formulas are also an explicit explanation of the cartesian
squares on page 047: restricting $p_{ij}$ to a base face is a
pullback. On page 044, with dimensions $r,r+s,r+s+t$, they give
$\dim U'=r+s-1$, $\dim V^{\circ\prime}=\dim U''=r+s+t-1$,
and $\dim\partial U''=r+s+t-2$.
Thus its item 6 must identify $\partial U''$ with
$\partial V^{\circ\prime}$ over $U$, not $\partial U$.
The original images of 039,043,044,047 were opened on 2026-11-19.
The omitted double prime in 044 is visible; the ambiguous subscript
in 047 is correct if read $i$ and is a dimensional error if read $j$.
\end{proof}

\begin{saresult}{A6.005}{Proved under stated hypotheses}{Reconstruction for any finite number of strata}
A resolution datum as in A6.004 reconstructs the compact stratified
space
\[
 X=\left(\coprod_{i\in I}M_i\right)\big/
 \bigl(x\sim p_{ij}(x),\ x\in H_{ij}\bigr).
\]
Its strata are exactly the images of the interiors $M_i^\circ$.
It has compatible conical control neighborhoods, with retraction
$\pi_j$ and radius $\rho_j$ induced by the boundary fibration and
its collar. Conversely, a compact finitely stratified space with
such conical control data has a resolution datum, obtained by
successively replacing cones on links by cylinders on the links,
starting at the closed strata. These operations are mutually inverse
when the collective faces, their common bases, and the collar data
are retained.

More precisely, the control data used in this statement include
locally trivial charts $V\times C(L)$ over each stratum, compact
links with the same structure recursively, and compatible radial
trivializations on overlapping tubes. Transition maps preserve
the cone rays and the radius, with the induced maps on links
compatible recursively. This is the conical-atlas condition in
Definition 1(v) of Albin--Leichtnam--Mazzeo--Piazza~\cite{ALMP2012},
\emph{The signature package on Witt spaces}, pp.249--250.
No Witt condition or orientability is needed in this topological
statement. The finite family formulation permits nonpure dimensions
and codimension-one strata.

The same constructions work in the PL category, and for finite
definable conical control atlases and compatible definable collars
in a fixed o-minimal expansion of a real closed field. In that
version compact means definably compact and the quotients carry
their definable-space topology. All chart and transition maps in
the construction stay in the given structure.
\end{saresult}
\begin{proof}
We give the quotient and the inverse explicitly, including the
induction needed for an arbitrary number of strata.

First, the relation has unique normal forms. If $x\in M_i$ lies on
the faces with lower labels $j$ and $k$, assume $j<k$. The two
first reductions are $p_{ij}(x)$ and $p_{ik}(x)$, and condition 2
says that the latter lies in $H_{kj}$ and reduces to the former.
Thus every two one-step reductions have a common reduction.
Every reduction strictly lowers the label in the finite poset,
so induction on the maximum remaining chain length proves that all
reduction sequences finish at the same point of some $M_j^\circ$.
This also proves that no distinct interior points are identified.

The relation is closed. Indeed, any sequence of reductions can
be replaced by a decreasing chain of length at most $|I|-1$.
For a fixed chain its domain is closed, its projection to its last
member is continuous, and the chain domain is compact.
Two points are equivalent precisely when some two such chains
have the same last value. This describes the relation as a finite
union of closed compact fiber products in
$(\coprod M_i)^2$. A closed equivalence relation on a compact
Hausdorff space has a compact Hausdorff quotient. To see the
separation assertion directly, two distinct classes are disjoint
compact sets. Choose disjoint open neighborhoods of them and
shrink each to the complement of the saturation of its closed
complement. Saturations of closed sets are closed, since they are
projections of closed subsets of the compact relation. The resulting
neighborhoods are saturated and still disjoint.

We prove the conical local assertion by induction on $|I|$.
It is immediate for a single manifold. Fix $y\in M_j^\circ$.
For each $i>j$ consider the fiber
\[
 K_i(y)=p_{ij}^{-1}(y).
\]
It is a compact manifold with corners. Its boundary faces are
exactly the $K_i(y)\cap H_{ik}$ with $j<k<i$, and restriction
of $p_{ik}$ maps them to $K_k(y)$.
All corner, fibration and collar identities restrict to these
fibers. They therefore form a resolution datum with fewer labels.
By induction its quotient is a compact stratified link $L_j(y)$.
By simultaneous local triviality, over a sufficiently small
$V\subset M_j^\circ$ this whole link construction is the product
of $V$ with that one link. The $j$-collars in every $M_i$ supply
one common radius coordinate $r_j\in[0,\epsilon)$. Their
compatibility says exactly that for $r_j>0$ the identifications
are those in $L_j(y)$, and at $r_j=0$ all points of a fiber are
identified with their base point in $V$. Hence the quotient
neighborhood is
\[
 V\times C_{\epsilon}(L_j(y)).
\]
The coordinate projection and radius define $\pi_j$ and $\rho_j$.
On an overlap with a higher tube $k>j$, the collar identities give
$\pi_j\pi_k=\pi_j$ and $\rho_j\pi_k=\rho_j$.
Incomparable tubes can be shrunk to be disjoint because their
face intersections are empty. The charts are smooth along strata,
and the product fibrations make $(\pi_j,\rho_j)$ submersive there.
This verifies the frontier, conical-chart and control requirements;
it is not only a dimension count.

For the inverse operation, retain any closed regular component
as its own boundaryless piece, and take a closed minimal singular
stratum $S_j$ in each remaining component.
Its control tube is the mapping cylinder of its compact link
bundle $L_j\to S_j$. Replace the tube by the corresponding
cylinder $L_j\times[0,\epsilon)$ over $S_j$, and glue to its
complement by the specified positive-radius trivialization.
In a chart this is the elementary replacement
\[
 V\times C_{\epsilon}(L)\quad\rightsquigarrow\quad
 V\times L\times[0,\epsilon).
\]
The new face is $L_j\to S_j$. Distinct minimal strata have
disjoint smaller tubes, so these replacements commute.
For another stratum, its inherited tube is obtained by doing
the identical replacement on its link and on its base when
necessary. The assumed radial compatibility means that these
descriptions agree on overlaps; in coordinates they are the same
projections with an additional nonnegative coordinate.
Every chain containing the removed minimal label is shorter by
one. Repeating the operation therefore terminates after at most
$|I|-1$ rounds. The remaining local models are products of
Euclidean open sets with finitely many half intervals. They are
manifolds with corners, and the closure of each lifted stratum
is $M_i$. A new face associated with $j<i$ projects to the
resolved closure $M_j$. Its fiber is the resolved link of $S_j$
in the closure of $S_i$. Local product charts give its fibration,
and two such faces meet precisely along the chains of lower
labels. The inherited projection identities are condition 2
of A6.004, and the inherited radius coordinates give condition 3.
They also prove the simultaneous local triviality used there.

On every control chart, replacing $C(L)$ by $L\times[0,\epsilon)$
and collapsing its zero-radius fibers are inverse operations.
The identifications on overlaps are the same transition maps.
They thus glue to global inverse identifications. This proves
mutual recovery for every finite poset, including several maximal
strata of different dimensions: the construction was performed
on each resolved stratum closure and glued through the specified
common lower bases, so it never assumed a unique top dimension.

All the preceding steps use finite products, closed fiber products,
the supplied local trivializations, and the explicit radial
replacement. They are PL when the input is PL, and definable
when the input is definable. For a definably compact realization,
the quotient relation is definable and definably closed; its
projections are closed on definable closed sets. Distinct classes
have disjoint definable neighborhoods by the minimum-distance
argument for disjoint definably compact sets. The saturation
argument above gives separated quotient charts. The finite
conical charts just constructed are therefore an atlas for the
definable quotient space. This avoids importing ordinary compactness
or differential-equation flows into a non-Archimedean field.

For comparison with the opened primary source, its Proposition 2.3
(pp.253--254) is the blowdown induction, Proposition 2.5
(pp.255--257) is the inverse resolution, and Remark 2.4 warns that
forgetting which disconnected hypersurfaces share a lower base
loses information. Our explicit collective-face labels retain
exactly that information. The simultaneous corner coordinates are
also the control data of its Definition 3 (p.252), and are developed
as iterated fibred tube systems in Debord--Lescure--Rochon~\cite{DLR2015},
Definitions 1.1 and 1.3 and Lemma 1.4 (pp.1804--1808).
\end{proof}

\begin{saresult}{A6.006}{Refuted}{What must be retained for categorical reconstruction}
For fixed resolutions, a stratum-preserving homeomorphism smooth
on each stratum and commuting with the scalar control functions
$\pi,\rho$ need not lift even continuously. This failure already
occurs for the compact semialgebraic suspension of a circle. Thus an equivalence
of categories requires a specified regularity condition on the
angular maps. One valid convention is that they extend smoothly
to the resolved corners, preserve the boundary fibrations, and
have radius components of the prescribed corner-map class.
With these morphisms and the complete data of A6.004--005,
resolution and blowdown are inverse functors. In the stricter
category preserving the chosen product trivializations, both
sides must use that same stricter condition. Forgetting the
common-base labels can even lose the
homeomorphism type of the reconstructed space.
\end{saresult}
\begin{proof}
Let $L=\mathbb{RP}^1$, and let $X$ be its suspension: in
$[0,1]\times L$ collapse each end circle to a separate point.
Use the two suspension points and the open cylinder as its three
strata. Its standard resolution is $[0,1]\times L$. For $0<r<1$ set
\[
 f([r,[u:v]])=[r,[ru:v]],
\]
and fix the two suspension points. For each $r>0$ the angular map is a projective linear
diffeomorphism, with inverse $[u:v]\mapsto[u:rv]$.
Both maps are semialgebraic and smooth on the open stratum.
They are continuous at the suspension points since they preserve
$r$. Near the lower point the control functions are the constant
retraction and $r$; near the upper point they are the constant
retraction and $1-r$. Thus both maps preserve all control functions.

Any lift to the fixed resolution must equal $(r,[u:v])\mapsto
(r,[ru:v])$ on $r>0$. Approach the boundary point $(0,[1:0])$
along $(r,[1:0])$ and along $(r,[1:r])$. The two image limits
are respectively $(0,[1:0])$ and $(0,[1:1])$. No continuous
extension exists. All hypotheses in the proposed weaker
morphism class have been checked, over any real closed field
using these two definable paths.

For the corrected morphism class, on $r>0$ a lift is prescribed
uniquely and the extension hypothesis supplies its value at
$r=0$. Iterating over corner coordinates proves existence of
the lift on every chart. On overlaps the lifts coincide on the
dense interior and hence everywhere, so they glue uniquely.
Conversely a compatible map of resolution systems is constant
on the equivalence classes in A6.005 and descends to a map of
quotients; its collar formulas give the stated conical
compatibility. These two operations respect composition and
are inverse on both objects and morphisms. This is the reason
for the lifted-map convention in Definition 4 of the opened
Albin--Leichtnam--Mazzeo--Piazza paper (p.258), rather than a
claim about every stratumwise smooth continuous map.

Finally, the same interval $M=[0,1]$ with its two boundary
points can be blown down by giving those points separate
point bases, producing an interval with two end strata, or
by grouping them as one collective face over one point,
producing a circle with one marked point. The separate
component fibrations are identical. The common-base grouping
distinguishes the results. This explains concretely the
nonuniqueness in the cited Remark 2.4.
\end{proof}

\begin{saresult}{A6.007}{Known}{Exactly what the peripheral groupoid of a compact surface determines}
For a compact surface $M$, the ordinary fundamental-groupoid functor
\[
 \Pi_1(\partial M)\longrightarrow\Pi_1(M),
\]
with its decomposition into connected components, determines the
unmarked homeomorphism type of every component of $M$, including
orientability, genus and number of boundary circles. For a connected
component with nonempty boundary, let $b$ be the number of boundary
circles and let $r$ be the rank of the abelianization of its
fundamental group. It is orientable precisely when
\[
 \ker\bigl(\mathbb Z^b=H_1(\partial M;\mathbb Z)
       \longrightarrow H_1(M;\mathbb Z)\bigr)\ne0.
\]
In the orientable case its genus is $(r-b+1)/2$; in the nonorientable
case its crosscap number is $r-b+1$. For a closed component,
$H_1=\mathbb Z^{2g}$ determines the orientable genus $g$, whereas
$H_1=\mathbb Z^{h-1}\oplus\mathbb Z/2$ determines the nonorientable
crosscap number $h$. Thus the sphere and the projective plane are
also distinguished.

This is the precise useful surface interpretation of the
``1-special'' and ``2-special'' ordinary groupoids on 121/043.
It does not assert that the manuscript defines those words:
the page contains a sketch and no definition of ``admissible
functor''. Corner flags and their attaching maps must still be kept
when several surface pieces are glued.
\end{saresult}
\begin{proof}
Choose one object in each connected groupoid. This recovers its
fundamental group, up to inner isomorphism, and the functor recovers
the homomorphisms from the boundary cyclic groups up to conjugacy.
Abelianization removes that conjugacy ambiguity, so the displayed
map is determined by the groupoid datum.

We derive the classification invariants using the opened
Gallier--Xu~\cite{GallierXu2025}, \emph{A Guide to the Classification Theorem for
Compact Surfaces}, Theorem 6.2, p.97, and the polygon presentations
of \S6.6, pp.101--103 (author version dated 8 January 2025).
For an orientable surface with $b>0$, the presentation is
\[
 \left\langle a_1,b_1,\ldots,a_g,b_g,d_1,\ldots,d_b
 \;\middle|\;\prod_{i=1}^g[a_i,b_i]\,d_1\cdots d_b=1\right\rangle.
\]
Eliminating $d_b$ gives a free group of rank $2g+b-1$.
In its abelianization the sole relation among boundary classes
is $d_1+\cdots+d_b=0$, up to the independent signs chosen for
boundary generators. The boundary kernel is rank one, generated
by a vector all of whose entries are $\pm1$.

For a nonorientable surface the presentation is
\[
 \left\langle a_1,\ldots,a_h,d_1,\ldots,d_b
 \;\middle|\;a_1^2\cdots a_h^2d_1\cdots d_b=1\right\rangle.
\]
Eliminating $d_b$ gives rank $h+b-1$. In the free abelianization
with basis $a_1,\ldots,a_h,d_1,\ldots,d_{b-1}$ one has
$d_b=-2\sum a_i-\sum_{j<b}d_j$. If $\sum n_jd_j=0$, the
coefficient of $a_1$ gives $n_b=0$, and then all other $n_j=0$.
The kernel is therefore zero. These calculations prove the
orientation test and both genus formulas. They also show that
in the orientable case the groupoid datum recovers the coherent
orientations of all boundary components up to simultaneous reversal.

If $b=0$, abelianizing the same presentations gives respectively
$\mathbb Z^{2g}$ and $\mathbb Z^{h-1}\oplus\mathbb Z/2$.
Theorem 6.2 states that orientability, number of boundary components
and Euler characteristic classify compact connected surfaces.
The recovered genera give the Euler characteristics
$2-2g-b$ and $2-h-b$, so the cited theorem applies and proves
the claimed complete recovery of the unmarked surface pieces.
The formulas also correct the isolated transposition of $p,q$
in the free-rank sentence on p.102 of the opened source:
the displayed presentation gives $2p+q-1$, as its abelianization
paragraph on that same page correctly states.

An interval's two end orientations form a $\mathbb Z/2$-torsor.
If the doubtful torsor on 043 were instead read $\mathbb Z/n$,
it cannot describe the end orientations of an edge unless $n=2$:
a torsor has exactly the cardinality of its acting finite group.
This settles both readings for that proposed geometric meaning.
\end{proof}

\begin{saresult}{A6.008}{Refuted}{Aggregate groupoids do not retain the corner gluing, even in dimension two}
Keeping only the peripheral groupoids of surface pieces, the
fundamental groupoid of the lower graph, and their attaching
functors does not determine a compact two-dimensional stratified
space up to stratified homeomorphism. The following two finite
polygonal gluings give an explicit counterexample:
\[
 w_1=aa^{-1}bb^{-1}cc^{-1},\qquad
 w_2=abb^{-1}a^{-1}cc^{-1}.
\]
In each case attach a closed hexagonal disk to the rose of three
circles by mapping each side linearly once around the edge named
by its signed letter. Use the rose vertex, its three open edges,
and the disk interior as strata. The aggregate groupoid diagrams
are equivalent, but the stratified spaces are not homeomorphic.
The finite program \texttt{code/check\_A6\_finite.py} checks both
words, their complete vertex links, and all signed edge relabelings
and polygon dihedral actions; its actual output is in
\texttt{logs/check\_A6\_finite.txt}.
\end{saresult}
\begin{proof}
Both spaces are compact finite polyhedra. Every point of an open
edge has the same two-sheet local model, and a sufficiently small
neighborhood of the vertex is the cone on its finite graph link.
They therefore have PL and semialgebraic conical control data;
the open disk and the three open edges are manifolds and the
frontier condition holds. Thus the examples satisfy the geometric
hypotheses, rather than merely giving abstract group presentations.

The disk groupoid is equivalent to the trivial groupoid, its
boundary groupoid to $B\mathbb Z$, and the rose groupoid to
$BF_3$. The boundary-to-disk functor is trivial. Both boundary
words freely reduce to the empty word, so both boundary-to-rose
functors are trivial as well. Each edge is used twice, with
opposite orientations. These facts identify the entire aggregate
diagram just specified, including each separate surface and graph
piece and its boundary functor.

Nevertheless the cyclic corner attachments differ. Name the six
half-edges of the rose $a,A,b,B,c,C$, with capital letters denoting
the opposite direction. A corner following a signed letter $x$
and preceding $y$ contributes a link edge from $-x$ to $y$.
For $w_1$ the link consists of a triangle on $a,b,c$ and one loop
at each of $A,B,C$. Thus the numbers of marked half-edge points
in its four connected components are $3,1,1,1$.
For $w_2$ there are two parallel edges between $A,b$, two
parallel edges between $a,c$, and loops at $B,C$. The corresponding
numbers are $2,2,1,1$.

This multiset is intrinsic: the components of a small punctured
vertex neighborhood are the components of the link, and the
marked points are the germs of the incident one-dimensional
strata. Any stratified homeomorphism preserves this incidence,
even if it permutes the three edge strata. The multisets differ,
so such a homeomorphism cannot exist. The program independently
obtains $[1,1,1,3]$ and $[1,1,2,2]$, and checks all
$3!2^3\cdot12=576$ signed relabeling/dihedral combinations.

The result identifies exactly the loss in this counterexample:
retaining the finite corner flags with their cyclic incidences
distinguishes the two gluings. It does not claim that the fuller
corner diagrams drawn above the groupoid sketch on 043 omit
those flags; the point is that replacing that fuller datum by
only the aggregate ordinary-groupoid functors is insufficient.
\end{proof}

\begin{saresult}{A6.009}{Proved}{The general information bound and the homotopy-colimit correction}
An ordinary fundamental groupoid, up to equivalence, determines
exactly the homotopy $1$-type: connected components and their
fundamental groups. A diagram of ordinary groupoids retains
the corresponding diagram of $1$-types, including its specified
functors and natural transformations. It cannot in general
recover higher homotopy within any piece, nor the homeomorphism
type of a stratified space. This remains false for compact
one-stratum smooth semialgebraic manifolds of fixed dimension:
$S^4$ and $\mathbb{CP}^2$ have identical fundamental groupoids.

In dimension two compactness and peripheral data are essential
to A6.007. The once-punctured torus and the thrice-punctured sphere
have equivalent fundamental groupoids and no manifold boundary,
but are not homeomorphic.

There is a precise positive homotopy statement. If a resolution
diagram consists of aspherical pieces, all attaching maps and
their coherence are retained, and its geometric realization is
the homotopy colimit of the diagram, then the groupoid diagram
determines that full homotopy type: take the homotopy colimit
of the classifying spaces of the groupoids. Taking a colimit
inside ordinary groupoids instead need not do so, even for a
sphere with one marked point.
\end{saresult}
\begin{proof}
A connected groupoid is equivalent to the one-object groupoid
of the automorphisms of a chosen object. Its classifying space
is $BG$. The nerve of the translation category of $G$ is
contractible (there is a unique arrow between every pair of
objects), and its free $G$-quotient is $BG$; hence $BG$ has
fundamental group $G$ and zero higher homotopy groups.
For a CW space $Z$, the natural map to the classifying space
of $\Pi_1 Z$ induces isomorphisms on $\pi_0,\pi_1$ and is
universal up to homotopy for maps from $Z$ to spaces with
vanishing homotopy groups above degree one. One direct proof
uses cellular obstruction theory: a homomorphism on $\pi_1$
defines the map on the two-skeleton, and all subsequent
extension and uniqueness obstructions lie in the zero groups
$\pi_n(BG)$, $n\ge2$. Applying this argument componentwise
proves the exact information bound. Functors and natural
transformations give the analogous maps and homotopies for
diagrams.

The sphere $S^4$ is simply connected and has $H_2(S^4;\mathbb Z)=0$.
The usual filtration
$\mathbb{CP}^0\subset\mathbb{CP}^1\subset\mathbb{CP}^2$
adds one open cell in each of dimensions $0,2,4$; its cellular
chain complex has one copy of $\mathbb Z$ in those degrees
and zero differentials. Its two-skeleton is $S^2$, and
attaching a four-cell does not change $\pi_1$, by cellular
approximation applied to loops and their homotopies.
Thus $\pi_1(\mathbb{CP}^2)=0$ but
$H_2(\mathbb{CP}^2;\mathbb Z)=\mathbb Z$.
Both are compact smooth four-manifolds with one stratum;
$\mathbb{CP}^2$ has a compact real algebraic realization as
the rank-one Hermitian projection matrices of trace one.
Their control data have no nontrivial boundary maps, and
their groupoids are both equivalent to the trivial groupoid.
Their different second homology proves that they are not
even homotopy equivalent.

For the noncompact surface example, both fundamental groups
are free of rank two by deleting a point from a torus or three
points from a sphere and retracting onto the usual graph
spines. The first surface has one end and the second has
three. To verify this without an appeal to an end-classification
theorem, take the compact complement of small disks about
the missing points. Every sufficiently large compact set
containing this complement leaves exactly one, respectively
three, nonrelatively-compact complementary components.
A homeomorphism preserves this number. Both surfaces are
smooth semialgebraic and have empty intrinsic manifold boundary.
They therefore give the same bare one-stratum groupoid data.

For the positive assertion, asphericity means that each natural
map $Z\to B\Pi_1 Z$ is a weak equivalence. Applied to the whole
coherent diagram, these give an objectwise weak equivalence.
The homotopy colimit is invariant under objectwise weak
equivalences: using its simplicial replacement, this follows
by replacing the diagram by a cofibrant one and successively
attaching the products with simplices along their boundaries.
Consequently its homotopy colimit is the claimed space.
The hypothesis that geometric gluing is a homotopy colimit
holds for the collared attachment diagrams here: replace each
attachment by its mapping cylinder. A collar identifies this
replacement with a deformation-equivalent geometric gluing;
iterate over the finite number of strata.

For the smallest illustrative case, the disk, its boundary,
and a point give the groupoid diagram
$1\leftarrow B\mathbb Z\to1$.
Its ordinary groupoid pushout is the trivial groupoid. The
homotopy pushout of its classifying spaces is
$*\cup^h_{S^1}*=\Sigma S^1=S^2$.
The map from the disk to the point is a homotopy equivalence,
and the collared boundary inclusion is a cofibration, so this
is exactly the homotopy type of $D^2/\partial D^2$.
This proves both the failure of ordinary groupoid colimits
and the corrected construction. Closed surface pieces such
as $S^2$ and $\mathbb{RP}^2$ must first be recovered by
A6.007, rather than replaced by their classifying spaces;
asphericity was an explicit hypothesis, not an automatic
consequence of being two-dimensional.
\end{proof}


\begin{saresult}{A6.010}{Proved under stated hypotheses}{A complete compact two-dimensional reconstruction datum}
For compact finite PL resolution systems of dimension at most two,
the following finite and groupoid data classify the blown-down
stratified spaces up to a homeomorphism preserving the resolution
labels:
\begin{enumerate}
\item the components of every resolved stratum, their dimensions,
and the finite corner flags, including the cyclic incidence of
arcs around each boundary circle and their common-base labels;
\item for each compact surface component, its peripheral functor
$\Pi_1(\partial M)\to\Pi_1(M)$, with the boundary components
identified with those in the flag datum;
\item for each one-dimensional boundary fibration over a resolved
one-dimensional stratum, its finite covering monodromy, with its
endpoint maps to the flags.
\end{enumerate}
Isomorphisms of these data must preserve the coherent diagrams,
including the boundary-loop maps in item 2. The one-dimensional
pieces themselves are compact intervals or circles, with their
finite endpoints; zero-dimensional pieces are finite sets.
For a trivial covering over an interval the covering datum is just
the set of sheets and the two endpoint maps. For a circle it is a
finite set with a permutation, up to simultaneous conjugacy and,
when the base orientation is reversed, inversion.

Thus ordinary peripheral groupoids determine the compact surface
pieces, and a complete flag/covering diagram supplies exactly the
remaining boundary assembly in dimension two. A6.008 proves that
one cannot discard the flag incidence and retain only aggregate
attaching groupoids. This statement does not classify mapping
spaces or replace the angular condition of A6.006.
\end{saresult}
\begin{proof}
Necessity follows by restriction of a homeomorphism respecting the
resolution diagram. For sufficiency we construct an isomorphism
of the resolved diagrams and apply A6.005. First match the finite
zero-dimensional pieces and the interval/circle components of the
one-dimensional pieces, preserving their endpoint flags. A
zero-dimensional proper local fibration over an interval is the
disjoint union of its sheets: lift the interval successively through
trivializing charts. Over a circle, lifting once round the circle
permutes a fiber, and the covering is recovered by gluing the two
ends of the trivial covering of an interval by that permutation.
This proves the stated complete covering classification and
constructs all required homeomorphisms of the faces lying over
one-dimensional bases. The faces over zero-dimensional bases
are compact one-manifolds with the given endpoint flags, so the
same interval/circle description matches them. The cyclic flags
make these matches into homeomorphisms of the entire surface
boundaries, compatible with all face fibrations and corner maps.

By A6.007 the peripheral functors determine the homeomorphism
types of the compact surfaces. We use the following elementary
boundary-extension fact: a homeomorphism of the boundaries of
two homeomorphic connected compact surfaces extends over the
surfaces if, in the orientable case, it either preserves all
coherently induced boundary orientations or reverses all of them;
in the nonorientable case every boundary homeomorphism extends.
Here boundary-component permutations are allowed.

To verify that fact, work with the standard disk-and-handle or
disk-and-crosscap model supplied by surface classification.
Any two boundary holes can be interchanged in a neighborhood
of an arc joining them; this is a pair of pants, and its two
inner holes can be exchanged while the outer boundary is fixed.
These exchanges realize every boundary permutation. An orientable
standard model admits a reflection reversing its orientation, so
one can realize either common boundary sign. In a nonorientable
model one may reverse just one prescribed boundary circle:
cap it by a disk, carry that disk once along a one-sided loop
avoiding the other boundary circles, and return it to its original
position. In the rectangular model of a one-sided band the return
identification is $(1,x)=(0,-x)$, so it reverses the disk. This
motion extends to an ambient homeomorphism supported in a small
neighborhood of the disk and its track: subdivide the track into
finitely many coordinate disks, move a smaller disk successively
inside each by a PL isotopy fixed at the chart boundary, and
compose these maps. Their supports avoid all other boundary
components. Removing the carried disk gives the claimed independent
boundary reversal. Repeating realizes every sign pattern.

After these permutations and signs have been matched, the remaining
map on each boundary circle preserves orientation. Parametrize the
circle by $\mathbb R/\mathbb Z$ and lift this map to an increasing
$h$ with $h(x+1)=h(x)+1$. The maps
$h_s(x)=(1-s)x+sh(x)$ are increasing homeomorphisms with the same
periodicity, and give an isotopy from the identity to the boundary
map. Apply this isotopy across a collar, with the identity at its
inner edge. The collars are disjoint, so these extensions combine.
This proves the boundary-extension fact with arbitrary prescribed
boundary maps, not just their isotopy classes.

For an orientable surface the necessary common-sign condition is
already enforced by the peripheral datum. By A6.007 the kernel
of $H_1(\partial M)\to H_1(M)$ is generated by the coherent sum
of the oriented boundary loops. An isomorphism of peripheral
diagrams carries this generator to itself or its negative; its
individual boundary signs are therefore all the same. The flag
and covering matches were required to induce precisely these
boundary-loop maps, so their assembled boundary homeomorphism
satisfies the extension condition. In the nonorientable case
there is no restriction. Extend over every surface component as
above; closed components need only the surface classification.
The resulting maps respect every fibration and every corner map,
hence descend to inverse homeomorphisms of the quotient systems.

This argument applies to finite semialgebraic or definable data
when the resolution and boundary maps have compatible finite PL
models: all coverings, finite gluing, collar isotopies and standard
surface constructions used in the proof can then be transported
through those models. It does not infer a classification of
noncompact surfaces from their ordinary groupoids; A6.009 explains
that explicit failure.
\end{proof}
